\section{Monte Carlo prediction of light collection and transport}

\subsection{Stage-resolved response}

The optical model should report separately the production of scintillation photons, WLS absorption, WLS re-emission, guided/transported photons, sensor arrivals, primary SiPM avalanches, correlated avalanches, charge and ADC response. Useful conditional efficiencies include
\[
\epsilon_{\rm WLS\,capture}=\frac{N_{\gamma,\rm WLS\,abs}}{N_{\gamma,\rm scint}},\quad
\epsilon_{\rm transport}=\frac{N_{\gamma,\rm sensor}}{N_{\gamma,\rm WLS\,emit}},\quad
\epsilon_{\rm PDE}=\frac{N_{\rm primary\,aval}}{N_{\gamma,\rm sensor}}.
\]
These ratios should be studied versus hit position, species/deposit regime and model nuisance parameters. Their denominators must be named explicitly: an efficiency conditional on WLS emission is not an absolute scintillator-to-sensor efficiency.

\subsection{Historical July calibration grid}

The July single-stave grid contains five proton/deuteron points and historically produced an order-ten detected-PE per visible-MeV scale. It is useful as a regression/history object, but it is not the nominal publication model.

\begin{publicationhold}
The exact July calibration revision predates multiple source-level fixes. In particular, inner and outer fibre claddings both used the shared NIST \texttt{G4\_PLEXIGLASS} material and attached different refractive-index tables, allowing the later assignment to overwrite the intended inner-clad optical properties. The campaign also predates explicit WLS fluorescence-yield/multiplicity, separated scintillator/Y-11 attenuation controls, corrected reflector-surface semantics, source-bound fibre-to-SiPM interface modelling and a single canonical SiPM response state. Its historical ``PE/MeV'' denominator is the Birks-visible variable, not unquenched deposited energy. Issue \#1303 therefore requires regeneration before any numerical light-yield result is promoted.

\subsection{Regenerated 2026-08 calibration grid (current model)}
Issue \#1303 regenerated the five-point grid (p60/p100/p140/d70/d110, 2000 events per point, seed-bound, distinct inner/outer cladding materials, SHA-256-bound inputs and manifests). The pooled Birks-visible light yield is $11.78~\mathrm{PE/MeV_{vis}}$ (offset $0.65$~PE, $r^2 = 0.976$, $n = 10{,}000$); per-point $E_{\rm vis}$ yields are flat within their bootstrap intervals (11.78--11.84), while $E_{\rm raw}$ yields (8.6--10.6) split by species through the Birks quenching ratio ($E_{\rm vis}/E_{\rm raw} \approx 0.81$ for 110~MeV deuterons). Stage accounting gives mean capture $\epsilon_{\rm WLS}\approx0.19$, transport $\approx0.021$, detection-given-arrival $\approx0.30$, and a saturated-PE-to-detected ratio $\approx0.96$ at the readout channel. These numbers are \textsc{mc\_model\_dependent}: the optical/SiPM nuisance envelope (PAPER-A07/A08) is not yet propagated, so the July grid remains the historical object and no beam-data calibration is claimed. Machine-readable source tables accompany every figure (\texttt{publication/tables/final/}).

A subsequent audit of the grid build revisions (both descendants of the
2026-08-11 change that introduced the \texttt{WLSMEANNUMBERPHOTONS} property)
shows the 2026-08 grid actually sampled \emph{Poisson}(1) WLS re-emission
multiplicity per absorption while its run metadata labelled the model
\texttt{geant4\_default\_one\_secondary}: in Geant4 11.2.2 the WLS process
samples Poisson whenever that property exists, even at mean 1. The pooled mean
light yield is unaffected in expectation ($E[\mathrm{Poisson}(1)]=1$), but
per-event multiplicity variance and saturation-adjacent observables carry
Poisson(1) semantics under a default-one label. Issue \#1088 now implements an
explicit three-mode multiplicity contract (exactly-one, Poisson($\mu$),
Bernoulli($q$) with $q=0.70$ from the Y-11/K27 quantum yield; Pla-Dalmau,
Foster, Zhang, NIM A361 (1995) 192--196), makes the optical digest
mode-complete, and adds a known-answer test that discriminates the modes; grid
regeneration under the corrected contract is tracked as a \#1303 follow-up.
\end{publicationhold}

The historical values and plots remain available under \texttt{publication/figures/gated/} and the source report for correction history only.

\subsection{Required nominal optical campaign}

A publication campaign must bind a clean source revision and build receipt; preserve \(E_{\rm raw}\) and \(E_{\rm vis}\) separately; record stage counters; use the declared one-fibre/one-end design channel as the nominal response only after the installed hardware configuration is source-bound under issue \#1296; retain alternative fibre/end observables as model controls where useful; scan or state the support in longitudinal/transverse position and angle; and propagate independent high-impact optical nuisance families rather than re-tuning the total mean response back to a historical target.
